% ==============================================================================
% Section 3: Theoretical Framework
% ==============================================================================
\section{Theoretical Framework}
\label{sec:theory}

\subsection{Cover Bidding in Procurement Auctions}

Consider a procurement auction for item $g$ at purchasing body $k$ in period $t$. Let $n$ denote the number of genuine (non-cartel) bidders and $m$ the number of cover bidders (FL firms). The procuring entity sets a reference price $r$ and requires a minimum of $\bar{n}$ bidders for the tender to proceed. If the number of participating bidders falls below $\bar{n}$, the tender is canceled and re-opened---potentially under a less competitive modality (e.g., direct purchase)---creating both delay costs for the procuring entity and uncertainty for the cartel.

A cartel with a designated winner faces a coordination problem: it must ensure that (i) the minimum bidder requirement is met, and (ii) the designated winner's bid appears competitive. Cover bidders solve both problems by submitting bids that are designed to lose. The cartel's problem is to choose $m$ (the number of cover bidders) and the cover bid distribution to maximize the designated winner's expected surplus, subject to the constraint that the auction appears legitimate to procurement officers. This trade-off generates two distinct equilibrium configurations, which we formalize as regimes.

\subsection{Two Regimes of Cover Bidding}

\begin{proposition}[Regime 1: Complementary Cover Bidding]
FL firms submit bids $b_j^{FL} = r + \epsilon_j$ where $\epsilon_j > 0$ is drawn independently to exceed the reference price. The designated winner bids $b^* = v + \delta$ where $\delta > 0$ is the cartel markup.
\end{proposition}

Under Regime 1, FL bids are high and dispersed because they are drawn independently above the reference price. The purpose is to satisfy the minimum bidder requirement ($n + m \geq \bar{n}$) without creating competitive pressure on the winner. This is the simpler coordination technology: cover bidders need only know the reference price $r$ and bid above it, with no need to observe the designated winner's bid or coordinate with other cover bidders. The cost to the cartel is transparency---an alert procurement officer may notice that several bidders submitted implausibly high bids.

\begin{proposition}[Regime 2: Coordinated Cover Bidding]
FL firms submit bids $b_j^{FL} = b^* + \eta_j$ where $\eta_j$ is small and positive, calibrated to create an illusion of competition around the winning bid.
\end{proposition}

Under Regime 2, FL bids cluster near the winning bid, reducing within-FL bid dispersion. The purpose is to make the auction outcome appear competitive---with several bids close to the winning price, procurement officers are less likely to suspect collusion. This is the more sophisticated coordination technology: it requires cover bidders to observe or be informed of the designated winner's bid $b^*$ and calibrate their bids accordingly. The cost is the risk of ``overshooting''---if a cover bidder's bid is accidentally lower than $b^*$, the cover bidder wins the contract, breaking the cartel arrangement.

\subsection{Comparative Statics}

The choice between Regime 1 and Regime 2 depends on the institutional environment. Regime 1 is less risky (no chance of accidental winning) but more detectable (high, dispersed bids). Regime 2 is harder to detect but requires more information exchange within the cartel. We expect Regime 1 to predominate in settings with (i) weak procurement oversight, (ii) many small purchasing bodies, and (iii) limited analytical capacity for detecting bid anomalies. These conditions characterize many sub-national procurement systems in developing countries, including the BEC platform.

\subsection{Testable Predictions}

\begin{prediction}[Price Effect]
\label{pred:price}
FL-present tenders have higher winning prices: $\beta_{\text{losers}} > 0$ in the price regression.
\end{prediction}

\begin{prediction}[Strategic Selection]
\label{pred:selection}
Cartels deploy cover bidders in tenders where the need to simulate competition is greatest. FL-present tenders therefore have weakly more total participants, and at least as many genuine competitors, as FL-absent tenders within the same item category.
\end{prediction}

\begin{prediction}[Regime Test]
\label{pred:regime}
Under Regime 1, $\text{Var}(b^{FL}) > \text{Var}(b^{\text{non-FL}})$. Under Regime 2, $\text{Var}(b^{FL}) < \text{Var}(b^{\text{non-FL}})$.
\end{prediction}

\begin{prediction}[Exchangeability Violation]
\label{pred:exchange}
FL bid residuals violate exchangeability (Bajari-Ye test), while non-FL bid residuals do not.
\end{prediction}

\begin{prediction}[Conditional Dependence]
\label{pred:independence}
FL bid residuals within the same tender are positively correlated. Non-FL bid residuals exhibit a weaker or zero correlation.
\end{prediction}

Predictions~\ref{pred:price}--\ref{pred:independence} are tested in Sections~\ref{sec:results}--\ref{sec:mechanisms}. We also explore the temporal dimension---whether FL entry into a market \textit{causes} price increases---through staggered DiD as a supplementary robustness exercise in Section~\ref{sec:robustness}. The hierarchy of predictions is deliberate: Predictions~\ref{pred:price}--\ref{pred:selection} establish the basic association, Prediction~\ref{pred:regime} distinguishes between organizational forms, and Predictions~\ref{pred:exchange}--\ref{pred:independence} provide direct statistical evidence of coordination through the Bajari-Ye framework.
